\chapter{Conclusiones y Trabajo Futuro}
\label{ch:conclusiones}

\section{Conclusiones}

Esta tesis se propuso comparar el rendimiento de los métodos de reconstrucción de canales EEG faltantes basados en Procesamiento de Señales en Grafos con regularización temporal frente a las referencias geométricas clásicas, utilizando un marco experimental riguroso, reproducible y exhaustivamente optimizado. Los resultados obtenidos permiten formular cinco conclusiones principales.

\textbf{Primera conclusión: la regularización temporal constituye un cambio de régimen, no una mejora incremental.} La evidencia experimental demuestra que la adición de un término de penalización temporal —incluso uno tan simple como la norma de la primera diferencia finita— transforma cualitativamente el comportamiento del interpolador. No se trata de una mejora marginal sobre los métodos espaciales, sino de un cambio de régimen: donde los métodos espaciales colapsan hacia ruido blanco bajo pérdida severa, TRSS preserva la estructura fina del espectro y la morfología de los componentes ERP con fidelidad clínica. El estudio de ablación confirma que este efecto proviene de una interacción sinérgica entre la información previa espacial y la temporal.

\textbf{Segunda conclusión: la implementación de MNE-Python es una referencia subestimada.} Lejos de ser un método obsoleto, la función de interpolación de MNE-Python es una referencia formidable gracias a seis mejoras de ingeniería acumuladas durante más de una década de desarrollo como código abierto, las cuales se documentan por primera vez en formato académico en esta tesis. En regímenes favorables, iguala o supera a TRSS en preservación espectral siendo entre 50 y 100 veces más rápido. Cualquier trabajo futuro que proponga un nuevo método de interpolación EEG debe compararse contra esta implementación, no contra versiones propias simplificadas.

\textbf{Tercera conclusión: la evaluación comparativa justa requiere optimizar todos los métodos.} Una contribución metodológica central de esta tesis es la demostración de que optimizar las referencias tanto como el método propuesto no es un lujo, sino una necesidad. La diferencia entre \textit{spherical splines} con parámetros por defecto y con parámetros optimizados mediante Optuna es sustancial —de 2 a 4~dB en SNR—. Sin esta optimización, cualquier método nuevo parecerá artificialmente superior.

\textbf{Cuarta conclusión: la elección del método debe ser condicional al contexto de aplicación.} Esta tesis rechaza explícitamente la narrativa de que TRSS es universalmente superior a MNE-Python. Los resultados muestran que la elección óptima depende del contexto: TRSS es superior cuando la preservación morfológica bajo pérdida severa es prioritaria; MNE-Python es superior cuando la velocidad de ejecución y la preservación espectral en regímenes de baja pérdida son prioritarios. Esta conclusión condicional es más útil para la práctica clínica que una declaración de superioridad universal, y es más defendible desde el punto de vista académico.

\textbf{Quinta conclusión: la trazabilidad y la reproducibilidad son factibles y necesarias.} Esta tesis demostró que es factible —con las herramientas adecuadas— diseñar una evaluación comparativa de interpolación EEG completamente trazable, donde cada número reportado puede rastrearse hasta el archivo de configuración y resultado que lo originó. El sistema de identificadores de ejecución, la separación estricta de fases exploratoria y confirmatoria, y el registro sistemático de versiones y semillas constituyen un estándar que debería adoptarse más ampliamente en la investigación en procesamiento de señales biomédicas.

\section{Trabajo Futuro}

Los resultados y limitaciones de esta tesis sugieren varias direcciones prometedoras para investigación futura.

En primer lugar, el término de penalización temporal empleado en TRSS —la primera diferencia finita— es el modelo más simple posible de continuidad temporal. Modelos más sofisticados, como información previa autoregresiva de orden superior que capture la estructura oscilatoria del EEG, o modelos de espacio de estados con dinámica de osciladores acoplados, podrían ofrecer mejoras adicionales, particularmente en la preservación de la estructura fina del espectro en bandas específicas.

En segundo lugar, el grafo de electrodos se construyó una vez y se mantuvo fijo para todo el registro. Sin embargo, la conectividad funcional del cerebro es dinámica: cambia con el estado cognitivo, la transición entre bandas de frecuencia y la presencia de eventos transitorios. Los grafos que se adaptan en el tiempo —reestimando los pesos en ventanas deslizantes basadas en coherencia instantánea— podrían capturar esta dinámica y mejorar la interpolación en estados no estacionarios.

En tercer lugar, este estudio midió la calidad de la interpolación en sí misma mediante métricas de señal. Una validación crucial pendiente es medir el impacto en tareas subsecuentes: ¿mejora TRSS la precisión de clasificación en interfaces cerebro-computador?, ¿reduce el error de localización de fuentes?, ¿preserva biomarcadores clínicos como la potencia relativa en bandas específicas para el diagnóstico de epilepsia o trastornos del sueño?

En cuarto lugar, el marco GSP con regularización temporal es conceptualmente aplicable a otras modalidades que comparten la estructura de señales muestreadas en dominios irregulares con dinámica temporal, como la magnetoencefalografía, la electrocorticografía y las matrices de microelectrodos.

En quinto lugar, la brecha de velocidad de entre 50 y 100 veces entre MNE y TRSS es un obstáculo práctico significativo. La investigación de implementaciones eficientes —aprovechando la estructura de banda del operador de diferencias temporales, precalculando descomposiciones del Laplaciano o utilizando resolvedores iterativos con convergencia temprana— podría reducir esta brecha sin sacrificar la calidad de reconstrucción.

Finalmente, la incorporación de TRSS como paso de interpolación dentro de los sistemas modernos de preprocesamiento automático de EEG, con selección automatizada de hiperparámetros basada en las características del registro, permitiría una evaluación en condiciones realistas de uso clínico.